% ARTIGO BASE
% https://arxiv.org/pdf/2302.02162

\documentclass[conference]{IEEEtran}
\IEEEoverridecommandlockouts
% The preceding line is only needed to identify funding in the first footnote. If that is unneeded, please comment it out.
\usepackage{cite}
\usepackage{amsmath,amssymb,amsfonts}
\usepackage{algorithmic}
\usepackage{graphicx}
\usepackage{textcomp}
\usepackage{xcolor}
\usepackage{flushend}
\def\BibTeX{{\rm B\kern-.05em{\sc i\kern-.025em b}\kern-.08em
    T\kern-.1667em\lower.7ex\hbox{E}\kern-.125emX}}

\usepackage{caption}
\usepackage{subcaption}

\usepackage[ruled,vlined,linesnumbered]{algorithm2e}

% partes para adequar ao camera ready
\pdfminorversion=5 % força PDF 1.5
\pdfobjcompresslevel=0 
\usepackage[left=1.62cm,right=1.62cm,top=1.9cm]{geometry}
% \usepackage[top=0.75in, bottom=0.75in, left=0.75in, right=0.75in]{geometry}

% original 0.90
\renewcommand{\baselinestretch}{1}
\setlength{\textfloatsep}{2.0pt plus 2.0pt minus 2.0pt}

\begin{document}


\title{Generic Title: Securing Semantic Communication}



\author{\IEEEauthorblockN{Prince~Nyarko\IEEEauthorrefmark{1},  Eirini~Eleni~Tsilopoulou\IEEEauthorrefmark{2}, 
André~Riker\IEEEauthorrefmark{1}\IEEEauthorrefmark{2}}
\IEEEauthorblockA{\IEEEauthorrefmark{1} Institute of Exact and Natural Sciences, Federal University of Par\'a, Brazil}
\IEEEauthorblockA{\IEEEauthorrefmark{2} Performance and Resource Optimization in Networks -- PROTON Lab\\ School of Electrical, Computer and Energy Engineering, Arizona State University, Tempe, AZ, USA}%%
}
%,castudillo@unicamp.br, weverton.cordeiro@inf.ufrgs.br, ariker@ufpa.br, gdecarvalho@brocku.ca
\maketitle

\begin{abstract}

Abstract ....

\end{abstract}

\begin{IEEEkeywords}
Federated Learning; 
\end{IEEEkeywords}

\section{Introduction}
% INTRO
Modern Machine Learning 

\section{Related Work}
\label{sec:related-work}

\subsection{Text Semantic Communication}
DeepSC formulates text transmission as a learned semantic communication task: Transformer-based processing supports sentence reconstruction over noisy channels, with the objective of recovering meaning rather than solely minimizing bit or symbol errors~\cite{xie2021deepsc}. Its use of sentence similarity alongside lexical reconstruction evaluation motivates assessing both semantic preservation and text fidelity~\cite{xie2021deepsc}. This framework provides the architectural context for studying knowledge transfer from a larger semantic communication Teacher to a compact Student.

\subsection{Efficient and Lightweight Semantic Communication}
Efficiency in semantic communication involves distinct model and transmission costs. L-DeepSC addresses deployment constraints through pruning and reduced weight precision, while also reducing the overhead of distributing model weights between the cloud/edge and devices~\cite{xie2021lite}. More recently, BidDeepSC-1.58b combines a shared bidirectional encoder--decoder architecture, low-bit quantization, and adaptable network widths~\cite{do2026biddeepsc}. These approaches illustrate that parameter storage and computation must be distinguished from the resources needed to transmit semantic representations. Accordingly, this paper treats Student model compression as its objective; reductions in channel use or bandwidth require separate evidence.

\subsection{Knowledge Distillation and Transformer Compression}
Knowledge distillation trains a Student using information supplied by a Teacher. The formulation of Hinton et al. uses softened output distributions to convey information beyond hard training labels~\cite{hinton2015distilling}. For sequence generation, Kim and Rush study word-level and sequence-level distillation in neural machine translation, demonstrating that the training targets and decoding procedure both matter when evaluating compact sequence models~\cite{kim2016sequencekd}. Their study concerns recurrent translation models rather than Transformer-based communication systems. TinyBERT provides a complementary Transformer-specific approach, transferring embedding, hidden-state, attention, and prediction information through general and task-specific distillation stages~\cite{jiao2020tinybert}. Together, these studies motivate examining the supervision available from a semantic communication Teacher, without prescribing which distillation losses the present implementation uses.

\subsection{Distillation in Semantic Communication and Research Positioning}
KD has already been applied directly to text semantic communication. Liu et al. study Transformer-based semantic encoders and decoders with neural channel coding under multi-user interference, combining knowledge transfer with model compression~\cite{liu2024kdsemcom}. Their evaluation includes compact models trained without KD, BLEU and sentence similarity, and varying noise and interference conditions~\cite{liu2024kdsemcom}. Thus, text modality, Transformer compression, and evaluation over noisy channels are already represented in the literature.

Other studies explore complementary deployment settings. Dynamic KD adapts semantic models to heterogeneous edge-device resource constraints~\cite{albaseer2024dynamickd}, while robust KD and lightweight architecture search transfer large-model capabilities to compact semantic encoders for image classification~\cite{ding2026robustkd}. A recent preprint by Eid et al. also studies compact DeepSC Students, combining teacher selection with residual channel representations and evaluating text quality and efficiency under channel variability~\cite{eid2026krumdeepsc}.

These precedents establish substantial overlap with the proposed setting~\cite{liu2024kdsemcom,eid2026krumdeepsc}. This work investigates the reconstruction-quality--model-complexity trade-off when a compact Transformer-based text semantic communication Student learns from a larger Teacher. Its intended evaluation separates the Teacher, a compact Student trained without KD where available, and the Student trained with KD under consistent channel and decoding conditions. The specific distinction from prior KD-based semantic communication remains provisional until the compression target, training objective, and experimental protocol are confirmed.

\section*{Acknowledgement}
This research was funded by CNPq (Brazil), under grants \#444978/2024-0, 201007/2025-8 and \#405026/2024-2.



\bibliographystyle{IEEEtran}
\bibliography{paper}

\end{document}